% A carefully factored condition-number example, not a timing claim.
\clearpage
\section*{Numerical audit: why a small residual can hide a large error}
\noindent Consider the nearly dependent matrix
\[
A_\delta=\begin{pmatrix}1&1\\1&1+\delta\end{pmatrix},
\quad
A_\delta x=b,\quad
x_*=\begin{pmatrix}2\\-1\end{pmatrix},
\quad
b=\begin{pmatrix}1\\1-\delta\end{pmatrix}.
\]
This exact problem has a unique answer whenever \(\delta\ne0\).
Yet uniqueness does not imply numerical stability: as \(\delta\)
shrinks, the columns become harder to distinguish from observations.

For \(\delta>0\), the trace and determinant of \(A_\delta^TA_\delta\) are
\[
\tau=4+2\delta+\delta^2,\qquad
\det(A_\delta^TA_\delta)=\delta^2.
\]
If its eigenvalues are \(\lambda_{\max}\) and \(\lambda_{\min}\),
their product is \(\delta^2\), and the singular values of \(A_\delta\)
are the square roots of those eigenvalues. Thus
\[
\kappa_2(A_\delta)
=\sqrt{\frac{\lambda_{\max}}{\lambda_{\min}}}
=\frac{\lambda_{\max}}{\delta}
\sim\frac4\delta,\qquad
\kappa_2(A_\delta^TA_\delta)=\kappa_2(A_\delta)^2.
\]
At \(\delta=10^{-3}\), the condition numbers are approximately
\(4.0020\times10^3\) and \(1.6016\times10^7\), respectively.
Normal-equation least squares inherits squared conditioning
\emph{before} accounting for additional floating-point rounding;
QR and SVD avoid that particular squaring in their solution pathways.

An observation-level illustration exposes the underlying information
problem, not just the algorithmic one. Hold \(b_1=1\), but perturb
\(b_2\) upward by \(10^{-4}\). Subtracting the two equations gives
\(\delta x_2=-\delta+10^{-4}\); for \(\delta=10^{-3}\),
the recovered coefficients move from \((2,-1)\) to
\((1.9,-0.9)\). A \(10^{-4}\) change in one observation
moves each coefficient by \(0.1\).
The new solution may have almost zero residual for the new data and
still be far from the true coefficients.

\paragraph{Implementation contract.}
When a dense least-squares problem is nearly rank deficient, report
\(\lVert Ax-b\rVert\), relative residual, rank or singular-value threshold,
the estimated condition number, and coefficient sensitivity.
Compare pivoted QR/SVD with normal equations using the \emph{same}
observations and numerical precision. A method can be backward stable
yet unable to recover information the data never constrained.

\paragraph{Challenge.}
Sweep \(\delta=10^{-1},10^{-2},\ldots,10^{-8}\).
Use a library reference implementation rather than promising that a
handwritten \(2\times2\) inverse is a stable general-purpose solver.
Plot condition number, observed coefficient error, and residual
on separate labeled axes. Discuss the scale at which floating-point
results cease to be trustworthy.


\begin{booktwolab}[Matrix conditioning]
\paragraph{C++23 checkpoint: a matrix problem with a numerical contract.}
Return to the chapter's conditioning discussion with the native
\texttt{book02\_workbench matrix} experiment. The exact example is
\[
 A=\begin{bmatrix}2&1\\1&3\end{bmatrix},\qquad
 x_\star=\begin{bmatrix}3/2\\-2\end{bmatrix},\qquad b=Ax_\star.
\]
Its answer is known before executing code, so this is a reconstruction
test rather than a comparison between two implementations that might
share a bug. The implementation scales the coefficient matrix, selects
the larger leading pivot and performs two-by-two elimination. It
rejects a pivot sufficiently small relative to its scaled matrix.
The rejection is a numerical policy, not an assertion of symbolic rank.

Now replace the second row of \(A\) with the first row multiplied by
two. The matrix is singular, so a solver must reject it; silently
returning huge finite components is unacceptable. Next consider
\[
 A_\varepsilon=
 \begin{bmatrix}1&1\\1&1+\varepsilon\end{bmatrix}.
\]
Derive its determinant \(\varepsilon\) and explain why the distance from
singularity is not the same thing as the error of a particular solution.
For fixed \(b\), the forward error depends on the conditioning of the
system as well as on roundoff and residual size. Compare the normalized
residual \(\lVert Ax-b\rVert/(\lVert A\rVert\lVert x\rVert+\lVert b\rVert)\)
with the forward error \(\lVert x-x_\star\rVert/\lVert x_\star\rVert\).
Neither quantity should be substituted for the other.

\end{booktwolab}

\paragraph{Independent check.}
Calculate the inverse of the benign \(2\times2\) matrix by hand, then
compare its prediction with the C++23 output. Perturb \(b\) and
\(\varepsilon\) separately to observe two different sensitivities.
Document the tolerance, the floating-point type, the exact matrix,
and the parameter range. Before claiming a general solver, implement
rank-revealing QR and independent higher-precision checks.


% Book Two next native C++23 pass: a rank-aware alternative to normal equations.
\begin{booktwoproof}[Householder QR instead of squaring the condition number]
For an overdetermined matrix \(A\in\mathbb{R}^{m\times n}\), \(m\ge n\),
the least-squares task is \(\min_x\|Ax-b\|_2\).
An orthogonal change of coordinates preserves Euclidean length:
if a product of Householder reflections yields
\(AP=Q\begin{bmatrix}R\\0\end{bmatrix}\), where \(P\) permutes
columns and \(Q^TQ=I\), then
\[
 \|Ax-b\|_2^2=
 \left\|R(P^Tx)-(Q^Tb)_{1:n}\right\|_2^2+
 \left\|(Q^Tb)_{n+1:m}\right\|_2^2 .
\]
The second term is independent of \(x\). Solving the upper
triangular system determines the first term without forming
\(A^TA\), whose condition number is ordinarily the square of
\(\kappa_2(A)\) for full-column-rank \(A\). Each reflection
\(H=I-2vv^T/(v^Tv)\) annihilates entries below a selected
diagonal. Column pivoting moves the largest remaining column
to the front, helping reveal loss of numerical rank.

The native \texttt{advanced::qr\_least\_squares} checks
input finiteness and uses a scale-relative diagonal threshold.
It reports a ratio of extreme \emph{QR diagonal magnitudes};
that ratio is a diagnostic, not a certified singular-value
condition estimate. At a detected deficient rank it rejects
the fit instead of inventing a unique coefficient vector.
\end{booktwoproof}

\begin{booktwolab}[QR sensitivity sweep in C++23]
The native command \texttt{book02\_workbench qr-sensitivity}
constructs a four-observation, two-parameter model with
rows \((1,1)\), \((1,1+\delta)\), \((1,1-\delta)\), and \((2,2)\).
Its generating parameter is \(x_\star=(2,-1)^T\).
Only the second observation is perturbed upward by \(10^{-10}\).
For \(\delta=10^{-1},\ldots,10^{-8}\), it exports
coefficient error against \(x_\star\), residual RMS, and the
QR diagonal diagnostic. All quantities come from C++23;
the optional Python utility only plots the recorded CSV.

Derive the exact unperturbed right-hand side
\((1,1-\delta,1+\delta,2)^T\). Predict why coefficient
error can grow as the columns approach dependence, even if
the residual remains tiny. Inspect the \emph{actual}
floating-point results rather than requiring an invented
smooth curve or monotonic progression. The supplied
implementation is a pedagogical dense pivoted QR,
not a complete LAPACK-style scaled rank-revealing algorithm.
\end{booktwolab}
